% Generated from the matching Typst scaffold by
% scripts/migrate_typst_report_to_latex.py (prose drafted 2026-06-25).
% INTEGRITY: model = deepseek-v4-flash. The RESULTS clause must NOT assert the pending
% headline -- it states the question; the finding is [RESULT: pending runs_144_seq_cceb325].
% No TOST claim. Thread "retrieval held constant", "non-inferiority", "footprint".

% ── Abstract ─────────────────────────────────────────────────────────────
% PURPOSE: 130-160-word structured abstract of the whole thesis.
% ─────────────────────────────────────────────────────────────────────────

\chapter*{Abstract}\addcontentsline{toc}{chapter}{Abstract}

LLM coding agents resolve real repository issues, but applying them across long
\emph{sequences} of related tasks is harder: experience accumulates as persistent memory, yet
naive accumulation inflates cost with no proven benefit and can even mislead. Rather than
invent another memory policy, we isolate \emph{storage} from \emph{retrieval}: six forgetting
policies, bracketed by No Memory and Full Memory, are compared under \emph{identical,
retrieval-held-constant} cosine search, a single frozen model (\texttt{deepseek-v4-flash}),
and a 144-run matrix over the eight SWE-Bench-CL sequences and three seeds. The analysis is
pre-registered and effect-size-first (Wilcoxon on $N=8$ sequence means, rank-biserial effect
sizes, bootstrap intervals), with cost read on two axes, compute and memory footprint. We ask
whether proactive forgetting is non-inferior to full accumulation, and at what footprint:
\texttt{[RESULT: pending runs\_144\_seq\_cceb325]}. The contribution is a controlled
\emph{measurement} --- including an instrument-validation sub-study that caught and fixed a
tooling defect before the matrix --- of \emph{when} and \emph{how much} a coding agent should
forget, not a new policy.
